% First post-TAC report, 2026-07-20.
% Compile from this directory so the relative figure paths resolve.
% Finalized meeting illustrations are stored in the adjacent figures directory.

\documentclass[11pt]{article}

\usepackage[margin=1in]{geometry}
\usepackage[T1]{fontenc}
\usepackage[utf8]{inputenc}
\usepackage{lmodern}
\usepackage{microtype}
\usepackage{amsmath,amssymb}
\usepackage{booktabs}
\usepackage{tabularx}
\usepackage{array}
\usepackage{graphicx}
\usepackage{xcolor}
\usepackage{caption}
\usepackage{float}
\usepackage{enumitem}
\usepackage{hyperref}
\usepackage{url}
\usepackage{fancyhdr}

\definecolor{navy}{HTML}{15324B}
\definecolor{teal}{HTML}{167D78}
\definecolor{orange}{HTML}{E76F51}
\definecolor{pale}{HTML}{F2F6FA}
\definecolor{muted}{HTML}{586B7D}

\hypersetup{
  colorlinks=true,
  linkcolor=navy,
  citecolor=teal,
  urlcolor=teal,
  pdftitle={Lib1 single-part sequence-to-expression baselines: July 2026 TAC report},
  pdfauthor={Minhang Xu}
}

\setlist{nosep,leftmargin=*}
\setlength{\parindent}{0pt}
\setlength{\parskip}{0.55em}
\renewcommand{\arraystretch}{1.16}
\pagestyle{fancy}
\fancyhf{}
\lhead{July 2026 TAC report}
\rhead{Lib1 mean-expression baseline}
\cfoot{\thepage}

\newcommand{\target}{\ensuremath{\log_2(\mathrm{total\ RNA}/\mathrm{total\ DNA})}}
\newcommand{\oof}{out-of-fold}
\newcommand{\cre}{cis-regulatory element}
\newcommand{\creplural}{cis-regulatory elements}
\newcommand{\code}[1]{\texttt{#1}}

% A compilation-safe figure macro. If the source is absent, the report retains
% a visible pointer rather than silently dropping the figure.
\newcommand{\assetfigure}[4]{%
\begin{figure}[H]
  \centering
  \IfFileExists{#1}{%
    \includegraphics[width=0.98\linewidth,height=0.78\textheight,keepaspectratio]{#1}%
  }{%
    \fcolorbox{navy}{pale}{%
      \parbox[c][0.25\textheight][c]{0.90\linewidth}{%
        \centering\textbf{Figure asset not found at compile time}\\[0.8em]
        \small\texttt{\detokenize{#1}}%
      }%
    }%
  }
  \caption[#2]{\textbf{#2} #4}
  \label{#3}
\end{figure}
}

\newcommand{\statusbox}[1]{%
  \begin{center}
  \fcolorbox{teal}{pale}{\parbox{0.92\linewidth}{#1}}
  \end{center}
}

\title{\textbf{Toward Objective-Driven Search of Combinatorial Gene-Regulatory Design Space}\\[0.35em]
\large Learning the sequence-to-expression landscape from Library 1}
\author{Minhang Xu}
\date{Thesis Advisory Committee meeting, 20 July 2026\\
Post-meeting report}

\begin{document}
\maketitle

\begin{abstract}
Mapping genotype to phenotype (the GP map) is a major effort in genomic
science. The problem is difficult because the search space expands from four
nucleotide choices at each position to combinations of entire regulatory
parts. In collaboration with the Bashor Lab at Rice University, we take a
synthetic-biology approach built around a defined DNA construct with five
cis-regulatory element (CRE) variable regions: Enhancer, Promoter, 5'UTR,
Intron, and 3'UTR. RNA and protein expression provide measurable phenotypes,
and the CLASSIC experimental framework makes it possible to connect these
phenotypes to large libraries of composed CRE circuits.

This report focuses on the Library 1 (Lib1) single-part libraries, in which one
CRE is varied at a time. The immediate objective is to establish a rigorous
sequence-to-expression baseline for all five parts using one construct-level
target, \target. The selected models show useful held-out predictive
performance across all five parts, although the best model and training choices
are part-specific. For Enhancer, models fit from random initialization were
substantially weaker during development; transfer learning produced the best
Enhancer model, although it still had the weakest final-test performance. Pooled Intron
performance is partly assisted by separation among sequence-defined groups.
The TAC discussion emphasized moving from individual parts to
multi-part composition models, comparing RNA and protein phenotypes, adding
simple k-mer and clustering baselines, and developing nucleotide- and
motif-level interpretation. Together, these directions move the project from
single-part prediction toward interpretable, model-guided design of regulatory
constructs.
\end{abstract}

\statusbox{\textbf{Report structure.} The report follows the July TAC
presentation from project scope and previous work through Lib1 preprocessing,
model development, and results. The later sections summarize the committee
discussion and next directions. Technical follow-ups developed after the
meeting are labeled in the corresponding results sections.}

\tableofcontents
\clearpage

\section{Project scope and thesis trajectory}

Genotype-to-phenotype mapping asks how changes in DNA sequence produce changes
in a measurable biological phenotype. Even before different regulatory parts
are combined, the nucleotide-level search space is enormous: a sequence of
length $L$ has $4^L$ possible designs. A genetic circuit adds a second layer of
combinatorial growth because prior-knowledge-defined regions can vary both
within themselves and in combination with one another.

Our collaboration with the Bashor Lab approaches this problem using a defined
synthetic DNA construct with five CRE variable regions---Enhancer, Promoter,
5'UTR, Intron, and 3'UTR---and RNA and protein expression as experimental
readouts. This scale of measurement is enabled by the CLASSIC experimental
framework (combining long- and short-range sequencing to investigate genetic
complexity), which links the full composition of multi-part constructs to
quantitative expression measurements through long- and short-read sequencing
\cite{rai2026ultra}.

The long-term goal is to learn a composition-to-function map that can guide the
search for constructs with a desired expression profile. This report begins
with the first controlled layer of that problem: the Lib1 single-part
libraries. Each sublibrary varies the sequence of one CRE while the other
regulatory regions are held at control baseline sequences. The present models
use mean RNA expression, \target, as their output.

Starting with individual CREs follows how the library is built and gives us a
reference for interpreting later multi-part behavior. The Lib1 baseline serves
three immediate roles:

\begin{enumerate}
  \item quantify how well expression can be predicted from each individual
        part using lightweight deep-learning models, with Pearson $r$ as the
        primary development metric;
  \item identify part-specific architecture and training choices before
        combining the parts; and
  \item estimate how performance changes with additional constructs and use
        those results to guide future library design.
\end{enumerate}

\assetfigure{figures/project_scope_composition_to_function_search.png}
{From individual CRE sequence models to composition-to-function search}
{fig:project-scope}
{A five-part regulatory construct creates both a nucleotide-level search space
within each CRE and a composition-level search space across CREs. CLASSIC
connects construct composition to measured expression.}

\section{Recap from the previous TAC}

The work leading into Lib1 had three parts. \textbf{First,} I used the published CLASSIC
data to test inference with sequence foundation models. That exercise showed
that a pretrained model such as Enformer can be applied to our construct
context, but its roughly 250-million-parameter scale makes full fine-tuning an
impractical starting point for the current data \cite{avsec2021enformer}.

\textbf{Second,} I used public single-CRE datasets to build and compare smaller CNN and
ResNet models. These models, ranging from roughly 300,000 to four million
parameters, are a better match for the shorter synthetic sequences in this
project. This work also established the CODA/BODA2-based software framework
used for hyperparameter optimization, training, and evaluation.

\textbf{Third,} the smaller Library 0 (Lib0) dataset made clear that a construct has not
only a mean expression value but also a spread across barcode observations.
Modeling that spread will require a probabilistic, count-based framework.
Lib0, however, was too small to support the sequence-level models that are the
focus of this report.

Lib1 is the first experimental dataset in the collaboration to contain a large
collection of single-part variants; it also contains multi-part sublibraries
that are not the focus of this meeting. The single-part data let us establish
the mean-expression baseline first, then compare it in future work with direct
inference from foundation models or lightweight predictors trained on
foundation-model embeddings, probabilistic observation models, and joint
RNA/protein outcomes.

\section{Library 1 data and construct-level target}

\subsection{Five single-part libraries}

Lib1 contains five single-part sublibraries. Within each sublibrary, one CRE
sequence is varied while the other regulatory regions are held at control
sequences. We train a separate model for each varied region, so each model
is built around one variable CRE. The from-scratch models receive that CRE
sequence directly, while the transferred Enhancer route places the variable
Enhancer inside a fixed 600-nt assay frame.

\begin{table}[H]
\centering
\caption{Lib1 single-part modeling sets.}
\label{tab:lib1sets}
\begin{tabularx}{\linewidth}{lXrr}
\toprule
Part & Sequence-length/input policy & All modeled constructs & $n_{\mathrm{barcode}}\geq8$ \\
\midrule
Enhancer & observed valid length 76--211 nt; neutral pad to 216 for
from-scratch training; selected transfer-learning route uses a 600-nt
assay-framed input & 4,787 & 1,229 \\
Promoter & observed valid length 41--51 nt; neutral pad to 51 for
from-scratch training & 7,893 & 1,931 \\
5'UTR & modal/exact 50-nt input & 8,331 & 1,797 \\
Intron & modal/exact 80-nt input & 7,848 & 1,326 \\
3'UTR & modal/exact 100-nt input & 6,845 & 775 \\
\bottomrule
\end{tabularx}
\vspace{0.4em}

\begin{minipage}{0.98\linewidth}
\footnotesize\textit{Note.} The observed input lengths differ from the rounded
schematic values used in Figure~\ref{fig:project-scope}. For example, the
modeled Enhancer sequences are not all exactly 200 nt, and observed Promoter
sequences are 41--51 nt rather than a fixed schematic length.
\end{minipage}
\end{table}

\subsection{From barcode observations to one target per construct}

Each barcode row contains a construct identity, a distinct barcode identity,
and paired DNA and RNA counts. Multiple barcode identities provide repeated
measurement support for the same construct.

For construct $i$ and barcode observation $j$, let $D_{ij}$ denote its DNA
count and $R_{ij}$ its RNA count. We aggregate the barcode observations as

\begin{equation}
D_i=\sum_j D_{ij}, \qquad
R_i=\sum_j R_{ij}, \qquad
y_i=\log_2\left(\frac{R_i}{D_i}\right).
\label{eq:target}
\end{equation}

where $D_i$ and $R_i$ are the aggregate DNA and RNA counts for construct $i$,
and $y_i$ is its construct-level expression target. This is a log ratio of
summed counts, not an arithmetic mean of per-barcode log ratios. No pseudocount
is added in this baseline, since every modeled construct has positive aggregate
DNA and RNA totals. During fitting, $y_i$ is z-scored using the current
training rows only. Predictions are inverse-transformed before reporting on the
raw \target\ scale.

\assetfigure{figures/lib1_barcode_observations_to_construct_target.png}
{From barcode observations to one expression target per construct}
{fig:barcode-to-target}
{Distinct barcode observations supporting the same construct are grouped by
variable-part sequence. DNA and RNA counts are summed before calculating one
\target\ response for that construct.}

\subsection{Expression and barcode-support distributions}

The raw expression distributions differ substantially across the five
sublibraries (Figure~\ref{fig:expression-distributions}). In the Promoter,
5'UTR, Intron, and 3'UTR sublibraries, the varied part is paired with a strong
control Enhancer, which keeps most constructs at relatively high RNA/DNA
expression. In the Enhancer sublibrary, that strong control is replaced by the
varied Enhancer sequence itself. The much lower and broader distribution shows
that varying the Enhancer has the most dramatic effect on RNA expression among
the five single-part libraries.

\assetfigure{../../../src/learn/outputs/analysis/lib1_dedup_data_summary_july2026/reporting/lib1_dedup_expression_target_distributions.png}
{Raw construct-expression distributions across five Lib1 single-part libraries}
{fig:expression-distributions}
{The panels show the raw construct-level target for all modeled constructs and
the higher-support ($n_{\mathrm{barcode}}\geq 8$) subset. The shift for the Enhancer library reflects removal
of the strong control Enhancer used in the other single-part libraries.}

Barcode support is likewise heterogeneous. The fraction of constructs with at
least eight distinct barcode identities is 25.7\% for Enhancer, 24.5\% for
Promoter, 21.6\% for 5'UTR, 16.9\% for Intron, and 11.3\% for 3'UTR.

\assetfigure{../../../src/learn/outputs/analysis/lib1_dedup_data_summary_july2026/reporting/lib1_dedup_barcode_support_distributions.png}
{Barcode-support distributions differ across the five Lib1 libraries}
{fig:barcode-support}
{The survival curves and composition bars describe distinct barcode identities
per construct.}

\section{Model development and evaluation workflow}

\subsection{Training set and higher-support evaluation set}

All constructs with at least one distinct barcode are eligible for model
training. Development and the final test use constructs with at least eight
barcodes. Earlier exact-bin and minimum-threshold
experiments showed no consistent advantage to discarding low-support training
constructs, while higher support provides a more stable population for model
comparison. A threshold of eight also preserved enough high-support 3'UTR
constructs for five development folds and a 250-construct final test.

Let $\mathcal{D}_{\mathrm{dev}}$ denote the development constructs and let
$f(i)$ be the fold containing construct $i$. For each fold, the model is trained
on the other four folds and then predicts the constructs in the fold that was
left out. Thus every development construct receives one prediction from a
model that was not trained on that construct:

\begin{equation}
\widehat y_i^{\mathrm{OOF}}=M_{-f(i)}(x_i),
\qquad
r_{\mathrm{pooled\ OOF}}=
\operatorname{PearsonCorr}_{i\in\mathcal{D}_{\mathrm{dev}}}
\left(y_i,\widehat y_i^{\mathrm{OOF}}\right).
\label{eq:oof-pearson}
\end{equation}

Here, $M_{-f(i)}$ is the model fit without fold $f(i)$, $x_i$ is the
variable-part sequence, $y_i$ is the observed construct-level expression target
from Equation~\ref{eq:target}, and $\widehat y_i^{\mathrm{OOF}}$ is the
corresponding held-out model prediction. The primary development score is one
Pearson correlation after the observed and held-out-predicted values from all
five folds are combined; it is not the average of five fold-specific
correlations.

One high-support final-test partition per part was fixed before model
development. After selection, the five part-specific model and training setups
were each refit at seeds 1701, 1702, and 1703, producing 15 trained model
checkpoints in total (five parts times three seeds). Here, a checkpoint is the
saved fitted model state for one part-specific setup and seed. For each part,
we averaged
the predictions from its three seed-specific models and evaluated that ensemble
once on the final test. We report those predictions without shifting or
rescaling them after observing the final-test results.

\subsection{Four stages of development}

The presentation organized the work as four controlled questions rather than
as an undifferentiated model search:

\begin{enumerate}
  \item Which architecture/optimization configurations are promising?
  \item Does RC augmentation help the same configuration?
  \item Does barcode-support weighting help the same arm?
  \item How does predictive performance change with the number of training
        constructs after model selection?
\end{enumerate}

\assetfigure{figures/lib1_single_part_modeling_workflow.png}
{Four-stage development workflow for the Lib1 single-part baselines}
{fig:modeling-workflow}
{The workflow holds the target and evaluation design fixed while comparing
model configurations, RC augmentation, barcode-weighted loss, and performance
as training-set size changes. Under the implemented Stage-3 formula, the
one-barcode weight is approximately 0.315; the 0.1 value shown in the current
illustration is the clipping floor and is not reached by retained constructs.}

\paragraph{Stage 1: broad configuration screen.}
We compared 885 configurations under a common fold-0, seed-1701, RC-off,
unweighted-MSE contract. The inventory comprised 128
Enhancer ResNet1D, 158 PromoterBassetVL, 156 Intron ResNet1D, 157 3'UTR
ResNet1D, 158 5'UTR ResNet1D, and 128 5'UTR UTRBassetVL configurations. Ten
configurations per part were promoted for the paired five-fold comparisons.

\paragraph{Stage 2: paired RC test.}
Each promoted configuration was evaluated with RC off and on while keeping
construct IDs, fold, seed, loss, and other training choices fixed. Validation
always used the original orientation. The completed product contained 660
development cells, 132 five-fold OOF arms, and 66 exact RC pairs. A limited
Enhancer-transfer lane and a limited 3'UTR UTRBassetVL lane were included.

\paragraph{Targeted 3'UTR follow-up.}
Because Stage 2 showed a promising but unstable 3'UTR UTRBassetVL region, a
fixed 24-configuration grid crossed learning rate, weight decay, and dropout
over five folds and two RC modes (240 cells). The result identified a useful
neighborhood around learning rate 0.002 and reinforced RC off for this part.
This targeted grid is not shown in Figure~\ref{fig:modeling-workflow} so that
the illustration can retain the common four-stage workflow.

\paragraph{Stage 3: paired barcode-weighted loss.}
For construct $i$ with $n_i$ distinct barcodes, let $z_i$ be the expression
target standardized using the current training rows. The weighted arm used

\begin{equation}
w_i=\operatorname{clip}\left[
\frac{\log(1+n_i)}{\log(9)},0.1,1\right], \qquad
L=\frac{\sum_i w_i(\widehat z_i-z_i)^2}{\sum_i w_i}.
\label{eq:weightedloss}
\end{equation}

The implemented weight is approximately 0.315 for a construct with one barcode
and reaches 1.0 at eight or more barcodes. Thus the 0.1 clipping floor is not
active for the retained constructs. Training contributions changed;
development and final metrics remained unweighted. Each weighted model was
compared with its exact unweighted mate. The final analysis contained 900 cells
and 180 complete five-fold OOF arms.

\paragraph{Stage 4: development-only downsampling.}
Selected policies were held fixed while training size varied over
$N=40,250,400,2500,4000,$ and full. Three frozen nested subset tracks were used
at each finite $N$, with five outer OOF folds and a distinct inner checkpoint
fold. All 660 cells completed. Final-test products were not read by Stage 4.

\section{Model-development results}

\subsection{Architecture and initialization route}

The first-order result was not a universal winning scratch architecture.
Enhancer required a distinct transfer route. With RC off on both sides, ten
scratch ResNet1D configurations had a median pooled OOF $r=0.113$ and maximum
$r=0.185$, whereas six transferred BassetBranched policies had median
$r=0.513$ and maximum $r=0.535$. The selected transferred K562/full policy
reached $r=0.565$ with RC on (Figure~\ref{fig:enhancer-transfer}).

This compares the two tested modeling routes; it does not isolate the effect of
pretraining because architecture, input framing, initialization, and source
training history also differ.

\assetfigure{../../../src/learn/outputs/analysis/lib1_dedup_tac_presentation_july2026/figures/main_enhancer_transfer_vs_scratch_oof.png}
{The tested transferred Enhancer route outperformed the tested scratch route}
{fig:enhancer-transfer}
{The route contrast is based on pooled five-fold OOF predictions. The smaller
within-transfer RC effect is shown separately from the much larger difference
between the from-scratch and transfer-learning routes.}

\subsection{RC augmentation is route specific}

Global RC augmentation consistently helped the transferred Enhancer family,
with a median mean-fold gain of approximately 0.031 across its six paired
policies. It was neutral or harmful on average for scratch Enhancer,
PromoterBassetVL, both UTR routes, and most Intron configurations. Only one
Intron configuration satisfied the stricter secondary checks. RC was
therefore selected only for the transferred Enhancer policy
(Figure~\ref{fig:rc-effect}).

My working hypothesis is that this route-specific benefit is related to the
training history of the transferred model. The BassetBranched checkpoint came
from the BODA2/Malinois training recipe, which exposed the source model to
reverse-complement sequences. Continuing that exposure during Lib1 fine-tuning
may fit the representation learned during pretraining better than it fits a
model initialized from scratch. This is a plausible explanation for the
pattern in Figure~\ref{fig:rc-effect}, not yet a causal test of why RC helps in
the specific Enhancer transfer-learning case.

\assetfigure{../../../src/learn/outputs/analysis/lib1_dedup_tac_presentation_july2026/figures/main_rc_augmentation_effect.png}
{Reverse-complement augmentation helped the transferred Enhancer route, not every CRE part}
{fig:rc-effect}
{Every comparison changes RC augmentation while holding the base
configuration, folds, seed, and loss fixed. The result supports a route-specific
training policy rather than a universal sequence symmetry.}

\paragraph{Follow-up to the RC result.}
The original BODA2 recipe combines three distinct choices: showing every
training sequence in both forward and reverse-complement orientation,
presenting high-expression sequences additional times, and averaging forward-
and reverse-strand predictions at inference \cite{gosai2024coda,boda2repo}.
The completed Lib1 comparison tests only the first choice globally.

A clearer follow-up begins without retraining: use the saved development
models to compare forward-only prediction with the average of forward and RC
predictions. A subsequent training comparison can give one arm an additional
RC copy and a matched control arm an additional forward copy of the same rows.
Both arms then see the same number of examples, so their difference more
directly tests orientation rather than repeated exposure. If that comparison
is informative, the same matched design can then test additional exposure for
high-expression or high-barcode-support constructs.

\subsection{Barcode-weighted loss gives modest, targeted gains}

Barcode weighting was selected for Promoter, Intron, 3'UTR, and 5'UTR, but not
Enhancer. The pooled OOF changes for the selected configuration in each part
were $-0.009$, $+0.009$, $+0.016$, $+0.010$, and $+0.023$, respectively.
Selection also required fold consistency and raw-error checks, so the
decision should not be reduced to a single pooled Pearson delta
(Figure~\ref{fig:weighted-effect}).

\assetfigure{../../../src/learn/outputs/analysis/lib1_dedup_tac_presentation_july2026/figures/main_barcode_weighted_loss_effect.png}
{Barcode-weighted loss produced small but repeatable gains in four selected configurations}
{fig:weighted-effect}
{Open points show held-out-fold changes and diamonds show the pooled OOF
change. Evaluation remained unweighted; the intervention changed only the
training loss.}

\subsection{Broad HPO landscape}

The configuration-level HPO landscape in Figure~\ref{fig:hpo-landscape} shows
that architecture and optimization choices were sampled broadly and that
several parts contain many similarly performing settings rather than one
isolated winning configuration. Each point uses the Stage-1 screening score:
validation Pearson $r$ on one predetermined split (fold 0, seed 1701), with RC
off and unweighted MSE. That single-fold score was used only to rank
configurations and promote candidates. The later decisions about RC
augmentation and barcode-weighted loss instead used matched five-fold
comparisons, testing each candidate against its paired RC-off/RC-on or
unweighted/weighted alternative. Those decisions considered pooled OOF
predictions together with consistency across folds and raw-error checks.

\assetfigure{../../../src/learn/outputs/analysis/lib1_dedup_tac_presentation_july2026/figures/supplement_hpo_configuration_landscape.png}
{Broad hyperparameter screening across five Lib1 single-part models}
{fig:hpo-landscape}
{The 885 tested settings span convolutional architecture, prediction heads,
regularization, optimizer, learning rate, weight decay, batch size, and
scheduler choices. Diamonds denote configurations promoted to paired five-fold
testing.}

The per-epoch histories answer a different question. As shown in
Figure~\ref{fig:training-dynamics}, the selected five-fold models generally
reach broad development plateaus rather than a uniquely sharp best epoch. This
current Lib1 figure is a more direct reference for training dynamics than the
earlier historical plot.

\assetfigure{../../../src/learn/outputs/analysis/lib1_dedup_tac_presentation_july2026/figures/supplement_selected_policy_training_dynamics_all_parts.png}
{Training dynamics of the selected five-fold policies}
{fig:training-dynamics}
{Training and held-out-fold trajectories are shown across epochs for the
selected policy in each CRE part. The histories are development-only and are
used to understand convergence and broad validation plateaus.}

\section{Selected models and held-out final test}

\begin{table}[H]
\centering
\caption{Development-selected policies and one-time final-test performance.}
\label{tab:selected-policies}
\small
\begin{tabularx}{\linewidth}{lXccrrrr}
\toprule
Part & Selected route & RC & Loss & OOF $r$ & Test $r$ & Test COD $R^2$ & Slope \\
\midrule
Enhancer & BassetBranched, K562/full transfer & on & unweighted & 0.565 & 0.365 & 0.004 & 0.504 \\
Promoter & PromoterBassetVL, scratch & off & weighted & 0.478 & 0.444 & 0.189 & 1.247 \\
Intron & ResNet1D, scratch & off & weighted & 0.690 & 0.681 & 0.462 & 0.950 \\
3'UTR & UTRBassetVL, scratch & off & weighted & 0.493 & 0.452 & 0.190 & 1.142 \\
5'UTR & UTRBassetVL, scratch & off & weighted & 0.542 & 0.512 & 0.257 & 0.927 \\
\bottomrule
\end{tabularx}
\end{table}

Four parts retained broadly similar association between development and the
one-time final test. Enhancer was the clear limitation: final-test Pearson fell
to 0.365, COD $R^2$ was approximately zero, and the observed-on-prediction
slope was 0.50. The model captures some rank ordering, but its raw predictions
do not match the observed expression scale well. Promoter and 3'UTR also show
departures from the identity line, whereas Intron and 5'UTR are closer
(Figure~\ref{fig:final-test}).

The plotted values are the predictions produced by the selected models. We did
not shift or rescale them after seeing the final-test results, so the slopes
describe the raw model outputs rather than a post-test-adjusted prediction.

\assetfigure{../../../src/learn/outputs/analysis/lib1_dedup_tac_presentation_july2026/figures/main_locked_final_test_scatter_hexbin.png}
{Locked final test: observed versus predicted construct expression}
{fig:final-test}
{The top row shows one point per construct and the bottom row summarizes the
same rows by hexagonal-bin counts. Dashed identity lines represent perfect raw-
scale calibration; the common solid fit color represents observed expression
regressed on prediction.}

\section{Understanding the Intron sublibrary effect}

The pooled Intron result is reproducibly strong, but it combines two sources of
signal: ranking constructs within sequence-defined groups and separating group
mean expression. Three nested IUPAC mask patterns define sequence-compatible
groups with markedly different target means. Because the masks overlap, these
assignments are sensitivity categories rather than verified synthesis-pool
labels.

On an illustrative Stage-1 held-out fold, the CNN achieved pooled
$r=0.778$. A leakage-safe predictor using only three group means fitted on the
training rows achieved $r=0.703$, while centering observations and predictions
within groups reduced the CNN correlation to $r=0.472$. This does not mean that
$0.778-0.472$ is an exact inflation estimate; the centered and pooled
correlations have different denominators. It does show why pooled Pearson alone
overstates uniform within-design-regime performance
(Figure~\ref{fig:intron-triptych}).

\assetfigure{../../../src/learn/outputs/analysis/lib1_dedup_tac_presentation_july2026/figures/main_intron_composition_triptych.png}
{Intron pooled performance contains substantial sequence-group separation}
{fig:intron-triptych}
{This is an explanatory Stage-1 fold diagnostic, not the selected final-policy
score. The middle panel uses group means fitted only on training rows. The
groups are inferred from sequence masks and are not measured splicing states.}

The result persists under the selected policy. Stage-2 OOF pooled Pearson was
0.682 versus 0.451 after within-group centering, with group-specific
correlations 0.599, 0.276, and 0.176. The locked final test produced pooled
0.681, centered 0.473, and group-specific 0.579, 0.338, and 0.206. Thus real
within-group signal remains, but the weakest residual group is not solved
(Figure~\ref{fig:intron-group-audit}).

\assetfigure{../../../src/learn/outputs/analysis/lib1_dedup_tac_presentation_july2026/figures/main_intron_sequence_group_audit.png}
{Selected-policy Intron performance within sequence-defined groups}
{fig:intron-group-audit}
{The selected model is summarized with pooled, within-group-centered, and
group-specific correlations for development OOF predictions and the held-out
final test. The groups are inferred from sequence masks rather than verified
experimental sublibrary labels.}

The Stage-2 variance decomposition in
Figure~\ref{fig:intron-group-audit} found that 70.6\% of target--prediction
covariance lay between inferred group means. Equalizing group weights barely
changed pooled Pearson because the group counts are already similar; it does
not remove the covariance created by different group means. Increasing the
barcode threshold also does not fix this composition problem.

The appropriate current claim is therefore:

\begin{quote}
The Intron model predicts the observed natural-mixture library well and learns
nonzero within-group sequence signal, but pooled performance is materially
assisted by sequence-group separation and performance is heterogeneous across
design regimes.
\end{quote}

\paragraph{Follow-up to the Intron result.}
The presentation described three Intron sublibraries, but the current analysis
can only infer three overlapping sequence groups from their design masks. The
next analysis will report performance within each group, request verified
synthesis-pool and design-family metadata from the Bashor Lab, and test a
fold-safe residual target that removes training-set group means before fitting
the sequence model. The pooled score remains useful for the observed mixture,
but within-group prediction is the more informative improvement target.

\section{In-silico downsampling: observed 10-fold and 100-fold gains}

Stage 4 holds each selected model and training setup fixed while repeating
development training with $N=40,250,400,2500,4000,$ and the full training set.
This creates directly observed tenfold comparisons within the downsampling grid
and one observed hundredfold comparison from 40 to 4,000 constructs. These
contrasts estimate how performance changes over the sampled training-size
range; they do not numerically forecast performance with ten or one hundred
times more data than the current full library.

\begin{table}[H]
\centering
\caption{Development-only Stage-4 gains across observed tenfold and hundredfold
training-size increases. Intervals are paired-bootstrap 95\% intervals.}
\label{tab:stage4}
\small
\begin{tabular}{lrr}
\toprule
Part/measure & $\Delta r$, 400 to 4,000 & $\Delta r$, 40 to 4,000 \\
\midrule
Enhancer & $+0.065\ [0.034,0.096]$ & $+0.128\ [0.076,0.183]$ \\
Promoter & $+0.194\ [0.148,0.242]$ & $+0.312\ [0.236,0.384]$ \\
Intron, pooled & $+0.078\ [0.043,0.117]$ & $+0.263\ [0.159,0.384]$ \\
Intron, within-group centered & $+0.231\ [0.153,0.313]$ & $+0.263\ [0.176,0.353]$ \\
3'UTR & $+0.242\ [0.070,0.389]$ & $+0.273\ [0.095,0.440]$ \\
5'UTR & $+0.254\ [0.196,0.308]$ & $+0.397\ [0.340,0.454]$ \\
\bottomrule
\end{tabular}
\end{table}

\assetfigure{../../../src/learn/outputs/analysis/lib1_dedup_stage4_downsampling_july2026/presentation/figures/01_primary_pearson_learning_curves.pdf}
{Predictive performance changes with the number of training constructs}
{fig:learning-curves}
{Finite-size points average three fixed, nested subset tracks. Error bars
preserve matched construct/subset comparisons through paired resampling. These
are development-only learning curves.}

Figure~\ref{fig:learning-curves} shows the full trajectories. The observed
400-to-4,000 gain is largest for 5'UTR, 3'UTR, and Promoter, while Enhancer and
pooled Intron change more modestly over the same decade. These statements
describe marginal gains over the sampled range rather than a ranking of which
CRE part should be prioritized experimentally.

The pooled Intron curve hides continued improvement within the inferred groups:
within-group-centered Intron performance increases by 0.231 from 400 to 4,000
constructs even though pooled performance increases by only 0.078. The paired
curves in Figure~\ref{fig:intron-learning-curves} make this distinction
explicit.

\assetfigure{../../../src/learn/outputs/analysis/lib1_dedup_stage4_downsampling_july2026/presentation/figures/03_intron_scoped_learning_curves.pdf}
{Intron learning curves for pooled and within-group prediction}
{fig:intron-learning-curves}
{The pooled curve measures performance on the observed mixture, whereas the
centered and group-specific curves track how much sequence variation is learned
within the inferred design groups.}

\section{Summary of Lib1 results presented at the TAC}

The results presented at the meeting support five working conclusions.

\begin{enumerate}
  \item \textbf{We established a common mean-expression baseline for all five
  individual CRE parts.} The data pipeline, target, development folds, and
  held-out evaluation are now defined consistently across the five
  sublibraries.
  \item \textbf{The best modeling choices are part-specific.} The selected
  Enhancer route uses a transferred BassetBranched model with RC augmentation.
  Promoter uses PromoterBassetVL, Intron uses ResNet1D, and both UTR models use
  UTRBassetVL; these four models are trained from scratch without RC.
  \item \textbf{All barcode-support levels are useful for training.}
  Higher-support constructs provide a more stable evaluation population, and
  barcode-weighted loss gives modest gains for four parts, but removing
  low-support constructs would discard useful sequence diversity.
  \item \textbf{The selected models generalize across all five CRE parts, but
  their limitations differ.} Enhancer has the weakest held-out performance.
  Intron has the strongest pooled correlation, but a substantial part of that
  score comes from separating sequence-defined groups.
  \item \textbf{More variants should improve the next generation of models.}
  The observed tenfold and hundredfold downsampling contrasts show continued
  gains as the number of training constructs increases.
\end{enumerate}

These results establish the single-part reference needed for the next stage;
they do not yet answer how the five CREs interact in a combined construct.
They also do not yet model protein output, expression spread, new cell types,
or an iterative design objective.

\section{TAC discussion and next directions}

\subsection{Move from individual CREs to multi-part composition}

The central next step is to use the multi-part Lib1 sublibraries to learn how
the five regulatory parts work together. The single-part models provide a
reference for asking whether a combined construct behaves as the sum of its
parts or contains higher-order interactions that require a joint model. We will
begin with transparent additive and modular baselines, then compare them with
models that learn directly from multiple part sequences. This is the first move
from sequence-to-function prediction for individual CREs toward the
composition-to-function search problem that motivates the thesis.

\subsection{Compare RNA and protein phenotypes}

RNA abundance is only one output of the synthetic construct. Because the Bashor
Lab measures protein expression as well, an immediate biological question is
whether different regulatory parts contribute differently to RNA and protein
output. From textbook gene regulation, we expect Promoter sequence to act most
directly on transcription, whereas the UTRs can strongly affect RNA stability
and translation. Comparing RNA with protein therefore gives us a way to test
whether the five parts contribute differently at transcriptional and
post-transcriptional stages. These measurements may also support models that
identify constructs with similar RNA but different protein expression, or
other deliberately separated RNA--protein profiles. This is a more immediate
design objective than cell-type-specific optimization, which will require
matched data from additional cell types.

\subsection{Add simple baselines and interpret sequence grammar}

The committee requested comparisons with models that are easier to inspect than
a CNN. We will add k-mer-based regression and sequence-clustering baselines
using the same folds, target, and evaluation metrics as the deep-learning
models. A motif-discovery analysis will be kept separate from the predictive
baseline so that discovered sequence patterns can be related back to expression
without overstating their causal role.

The committee also asked whether the models can be interpreted at nucleotide
and motif resolution. One potential plan is a staged adaptation of BODA2
interpretation methods \cite{gosai2024coda}: first build a part-specific,
single-output model adapter; then compare in-silico mutagenesis with Sampled
Integrated Gradients on a small development subset; validate high-contribution
regions through matched sequence perturbations; and only then scale to
TF-MoDISco Lite and motif annotation
\cite{sundararajan2017ig,shrikumar2018tfmodisco}. Intron interpretation will be
stratified by sequence-defined group so that group identity and splice
boundaries are not mistaken for fine-grained regulatory grammar.

Evolutionary conservation and disease-associated variants were raised as
possible downstream connections for interpretable sequence features. The
current synthetic constructs are deliberately different from naturally evolved
genomic loci: their part lengths and arrangement are engineered, and the
current assay uses plasmids rather than constructs integrated into the genome.
We therefore should not assume a direct conservation or disease interpretation.
We can nevertheless compare each synthetic CRE part with analogous endogenous
regulatory regions and publicly curated CRE profiles, examine whether the same
nucleotide- or motif-level features are important, and synthesize matched
variants inspired by known regulatory alleles to test their effects within the
controlled Lib1 construct context.

\subsection{Use foundation-model representations pragmatically}

Full fine-tuning of a large sequence foundation model is not the immediate
priority. A more feasible comparison is to extract frozen representations from
models such as Evo 2 and train a lightweight predictor on top. This keeps the
Lib1 task and evaluation splits unchanged while asking whether a pretrained
sequence representation adds information beyond one-hot CNN inputs.

\assetfigure{figures/evo2_embedding_prediction_plan.png}
{Potential use of Evo 2 embeddings for Lib1 expression prediction}
{fig:evo2-plan}
{A frozen foundation model supplies candidate intermediate-layer
representations, and a lightweight supervised head maps those representations
to expression phenotypes. This is a conceptual plan: the exact construct
length, token-position pooling, and RNA/protein target transformations remain
to be finalized.}

\subsection{Connect the project to regulatory-grammar literature and manuscript framing}

The committee recommended a closer review of MPRA and regulatory-grammar work
from Jay Shendure's group, together with Martin Kircher's work on experimental
design and measurement variability. That literature will help position what is
new about varying multiple regulatory elements in one synthetic construct.

The meeting also raised manuscript strategy. The main biological story is the
transition from individual CREs to multi-part regulatory grammar. In parallel,
some computational components may support focused methods or bioinformatics
papers. The relevant novelty is not that each machine-learning technique is
new in isolation, but that the techniques are combined, tested, and adapted for
this biological design space. We will make that distinction explicitly as the
work develops.

\subsection{Longer-term thesis trajectory}

After the mean-expression and multi-part baselines are stable, a probabilistic
observation model can return to the barcode-level RNA/DNA counts and model the
spread of expression in addition to its mean. Additional cell-type data would
then make cell-type-specific objectives possible. Ultimately, the models should
enter an iterative design--build--test--learn cycle that proposes constructs
with desired phenotypic profiles rather than serving only as retrospective
predictors.

For future meetings, the written TAC report will be circulated approximately
one week in advance and will retain enough project background to orient newer
committee members.

\section*{Closing perspective}

The July 2026 TAC marks a transition from establishing single-part baselines to
using them as a reference for composition-level questions. We now have a common
construct-level RNA target, part-specific model and training choices, and an
in-silico estimate of how performance changes as more training constructs are
added. The Enhancer and Intron results also identify concrete technical
questions that should be resolved as the baseline develops.

The central next phase is biological: combine individual CRE models, measure
higher-order interactions, and compare how Promoter, UTR, Enhancer, and Intron
variation contributes to RNA and protein phenotypes. Simple k-mer and
clustering baselines, sequence-level interpretation, and foundation-model
embeddings will help determine what the deep models add. Together, these steps
move the project toward composition-to-function prediction and model-guided
design of regulatory circuits.

\begin{thebibliography}{99}

\bibitem{rai2026ultra}
Rai, K. et al.
Ultra-high-throughput mapping of genetic design space.
\textit{Nature} \textbf{650}, 1035--1044 (2026).
\url{https://doi.org/10.1038/s41586-025-09933-9}.

\bibitem{gosai2024coda}
Gosai, S. J. et al.
Machine-guided design of cell-type-targeting cis-regulatory elements.
\textit{Nature} \textbf{634}, 1211--1220 (2024).
\url{https://doi.org/10.1038/s41586-024-08070-z}.

\bibitem{boda2repo}
Gosai, S. J. et al.
\textit{boda2: Computational Optimization of DNA Activity (CODA)}.
\url{https://github.com/sjgosai/boda2}.

\bibitem{avsec2021enformer}
Avsec, Z. et al.
Effective gene expression prediction from sequence by integrating long-range
interactions.
\textit{Nature Methods} \textbf{18}, 1196--1203 (2021).
\url{https://doi.org/10.1038/s41592-021-01252-x}.

\bibitem{kelley2016basset}
Kelley, D. R., Snoek, J. and Rinn, J. L.
Basset: learning the regulatory code of the accessible genome with deep
convolutional neural networks.
\textit{Genome Research} \textbf{26}, 990--999 (2016).
\url{https://doi.org/10.1101/gr.200535.115}.

\bibitem{sundararajan2017ig}
Sundararajan, M., Taly, A. and Yan, Q.
Axiomatic attribution for deep networks.
In \textit{Proceedings of the 34th International Conference on Machine
Learning} (2017).
\url{https://proceedings.mlr.press/v70/sundararajan17a.html}.

\bibitem{shrikumar2018tfmodisco}
Shrikumar, A. et al.
Technical note on transcription factor motif discovery from importance scores
(TF-MoDISco).
\textit{arXiv:1811.00416} (2018).
\url{https://arxiv.org/abs/1811.00416}.

\bibitem{bailey2021streme}
Bailey, T. L.
STREME: accurate and versatile sequence motif discovery.
\textit{Bioinformatics} \textbf{37}, 2834--2840 (2021).
\url{https://doi.org/10.1093/bioinformatics/btab203}.

\bibitem{sinai2020adalead}
Sinai, S. et al.
AdaLead: a simple and robust adaptive greedy search algorithm for sequence
design.
\textit{arXiv:2010.02141} (2020).
\url{https://arxiv.org/abs/2010.02141}.

\end{thebibliography}

\end{document}
